\documentclass[11pt]{article}

\usepackage[margin=1in]{geometry}
\usepackage[T1]{fontenc}
\usepackage[utf8]{inputenc}
\usepackage{enumitem}
\usepackage{hyperref}
\usepackage{lmodern}
\usepackage{setspace}
\usepackage{xcolor}
\usepackage{verbatim}

\hypersetup{
  colorlinks=true,
  linkcolor=blue,
  urlcolor=blue
}

\setlength{\parskip}{0.4em}
\setlength{\parindent}{0pt}

\begin{document}

\begin{center}
{\LARGE \textbf{Assurance Case Outline (ACO) Specification v1.0}}\\[0.5em]
\today
\end{center}

A human-friendly, indentation-based assurance case outline format designed for authors, supporting automatic transformation into an internal Node Model, APL (Assurance Pattern Language) blocks, and graphical representations.

\bigskip
\hrule
\bigskip

\section{Purpose and Scope}

This specification defines a structured textual outline language (``Assurance Case Outline'' -- ACO) for authoring assurance case patterns and instances.

It is:
\begin{itemize}[leftmargin=2em]
  \item Initially intended for use in the MedSecurance project.
  \item Nevertheless domain-independent and platform-neutral.
  \item Not tied to O-ETB internals, but consistent with O-ETB's APL concepts and with the GSN Community Standard where that does not conflict with the current O-ETB implementation.
\end{itemize}

The ACO notation aims to balance:
\begin{itemize}[leftmargin=2em]
  \item Ease of authoring (minimal punctuation, close to ``normal'' outlines).
  \item Readability (GSN-like, hierarchical structure).
  \item Machine-interpretability (convertible to APL, DOT/Graphviz, and other formats).
\end{itemize}

The \textbf{ACO outline is the human-facing source of truth}; APL and diagrams are generated from it.

In particular:
\begin{itemize}[leftmargin=2em]
  \item \textbf{ACO is unambiguous enough} to be parsed into a precise Node Model.
  \item The Node Model can then be systematically mapped to APL blocks and graphical representations.
  \item The Node Model is also the place where checks, diagnostics, and transformations are performed.
\end{itemize}

\bigskip
\hrule
\bigskip

\section{Core Concepts}

\subsection{Node Types}

ACO adopts a minimal, GSN-inspired set of node types, aligned with APL concepts:

\begin{itemize}[leftmargin=2em]
  \item \textbf{Goal} --- a claim that must hold.
  \item \textbf{Strategy} --- a structure explaining how children collectively support a parent (may be optional in early usage).
  \item \textbf{Context} --- assumptions, definitions, scope, environmental conditions.
  \item \textbf{Assumption} --- asserted conditions that are taken as given.
  \item \textbf{Justification} --- rationale explaining why an approach is acceptable.
  \item \textbf{Module} --- a reusable or composite block, such as:
  \begin{itemize}[leftmargin=1.5em]
    \item a pattern module,
    \item a component-level argument,
    \item an evidence-aggregation module.
  \end{itemize}
\end{itemize}

Additional types (e.g.\ Evidence blocks or other specialised constructs) can be introduced later, but the initial ACO v1.0 focuses on these.

\subsection{Node Header Syntax}

Each node is introduced by a \textbf{header line}, followed by one or more indented \textbf{body lines}.

The header has the form:
\begin{quote}
\texttt{\textless Type\textgreater{} [ID] \textless shortLabel\textgreater:}
\end{quote}

where:
\begin{itemize}[leftmargin=2em]
  \item \texttt{\textless Type\textgreater{}} is one of:
  \begin{itemize}[leftmargin=1.5em]
    \item \texttt{Goal}
    \item \texttt{Strategy}
    \item \texttt{Context}
    \item \texttt{Assumption}
    \item \texttt{Justification}
    \item \texttt{Module}
  \end{itemize}

  \item \texttt{[ID]} is an \textbf{optional} identifier, such as \texttt{G0}, \texttt{G1}, \texttt{C3}, \texttt{M2}, etc.
  \begin{itemize}[leftmargin=1.5em]
    \item If present, it is a single token (no spaces).
    \item If absent, the header consists only of \texttt{\textless Type\textgreater{} \textless shortLabel\textgreater:} (see below).
  \end{itemize}

  \item \texttt{\textless shortLabel\textgreater{}} is a single-token name (no spaces) describing the node.
  \begin{itemize}[leftmargin=1.5em]
    \item CamelCase or snake\_case is recommended (e.g.\ \texttt{operationalPlane}, \texttt{componentBehaviour}, \texttt{plane\_definition}).
  \end{itemize}

  \item The header line \textbf{must end} with a colon (\texttt{:}).
\end{itemize}

Thus:

\textbf{With ID} (three tokens: Type + ID + label):
\begin{verbatim}
Goal G0 operationalPlane:
  The operational plane guarantees the {local policy} is met.
\end{verbatim}

\textbf{Without ID} (two tokens: Type + label):
\begin{verbatim}
Context planeDefinition:
  The system model defines the plane and its properties.
\end{verbatim}

Both forms are allowed; the parsing rules treat the header uniformly and recover Type, optional ID, and label deterministically.

Parsing conventions (informal):
\begin{itemize}[leftmargin=2em]
  \item The first token is always the \textbf{Type}.
  \item If the second token matches the pattern of an ID (e.g.\ \texttt{G\textbackslash d+}, \texttt{C\textbackslash d+}, \texttt{M\textbackslash d+}, etc.), it is treated as the \textbf{ID} and the next token is the \textbf{label}.
  \item Otherwise, the second token is treated as the \textbf{label} and there is no ID.
  \item Everything between the label and the trailing colon \texttt{:} must be whitespace only.
\end{itemize}

\subsection{Node Body Text}

The header line is followed by one or more \textbf{indented body lines}, containing the natural-language description.

Example:
\begin{verbatim}
Goal G0 operationalPlane:
  The operational plane guarantees the {local policy} is met.

Context C0 planeDefinition:
  The system model describes the plane and its {properties}.

Module M1 component:
  Behaviour of the {component} meets its local policy.
\end{verbatim}

Rules:
\begin{itemize}[leftmargin=2em]
  \item Body lines must be indented \textbf{more} than the header (at least one indentation level deeper).
  \item Body text may span multiple lines.
  \item Blank lines are allowed between nodes for readability, but should not break the structural semantics.
\end{itemize}

\bigskip
\hrule
\bigskip

\section{Indentation-Based Structure}

\subsection{Indentation Rules}

Indentation is the primary way to represent structure and relationships.

\begin{itemize}[leftmargin=2em]
  \item Indentation uses \textbf{2 spaces per level}.
  \begin{itemize}[leftmargin=1.5em]
    \item Tabs are discouraged; tools should normalise or flag tabs.
  \end{itemize}

  \item A node header at indentation level \(k\) (\(k \geq 0\)):
  \begin{itemize}[leftmargin=1.5em]
    \item If \(k = 0\), it is a top-level node.
    \item If \(k > 0\), it is structurally a \textbf{child} of the nearest preceding node header at indentation level \(k - 1\).
  \end{itemize}

  \item When the indentation decreases from level \(k\) to level \(j\):
  \begin{itemize}[leftmargin=1.5em]
    \item All nodes at levels \(> j\) are considered ``closed''.
    \item The new node at level \(j\) is attached (as child or sibling) appropriately in the tree implied by indentation.
  \end{itemize}

  \item Indentation must never ``jump'' by more than one level without passing through the intermediate levels. For example:
  \begin{itemize}[leftmargin=1.5em]
    \item Going directly from level 0 to level 3 is \textbf{invalid}.
    \item This is treated as a structural error (see Section~\ref{sec:errors}).
  \end{itemize}
\end{itemize}

\subsection{Default Relationship Semantics}

From the indentation-defined tree, ACO derives \textbf{default relationships}.

For a parent node \(P\) and a child node \(N\) (implied by indentation):
\begin{itemize}[leftmargin=2em]
  \item If \(N\) is of type Goal, Strategy, Module, Assumption, or Justification:
  \begin{itemize}[leftmargin=1.5em]
    \item A default edge \texttt{supported\_by(P, N)} is added to the Node Model.
  \end{itemize}

  \item If \(N\) is of type Context:
  \begin{itemize}[leftmargin=1.5em]
    \item A default edge \texttt{in\_context\_of(P, N)} is added.
  \end{itemize}
\end{itemize}

This means:
\begin{itemize}[leftmargin=2em]
  \item Authors do \textbf{not} need to explicitly spell out relations for the common tree-structured argument.
  \item Non-tree or cross-cutting relations (DAG edges) can be added explicitly using natural-language relation sentences (see Section~\ref{sec:nl-relations}).
\end{itemize}

\bigskip
\hrule
\bigskip

\section{Optional Natural-Language Relations}
\label{sec:nl-relations}

To improve readability and to express \textbf{non-hierarchical} (non-tree) relations, authors may add optional explicit relation sentences.

These are written in an NL-ish style but parsed into relations over node IDs.

\subsection{Supported Relations}

The following patterns are accepted (case-insensitive, minor variations allowed).

\textbf{Supported-by:}
\begin{verbatim}
G0 is supported by G1, G2, C0 and M0.
G1, G2, C0 and M0 support G0.
\end{verbatim}

\textbf{In-context-of:}
\begin{verbatim}
G0 is in context of C0.
C0 provides context for G0.
\end{verbatim}

Where:
\begin{itemize}[leftmargin=2em]
  \item Left-hand side and right-hand side refer to existing node IDs (e.g.\ \texttt{G0}, \texttt{C0}, \texttt{M0}).
  \item The same IDs must appear as node IDs in headers.
\end{itemize}

\subsection{Semantics and Interaction with Indentation}

\begin{itemize}[leftmargin=2em]
  \item These explicit relation sentences \textbf{augment} the default relations implied by indentation.
  \item They may:
  \begin{itemize}[leftmargin=1.5em]
    \item reinforce links already implied by the hierarchy, or
    \item introduce additional DAG edges (e.g.\ supporting relations that cut across branches of the tree).
  \end{itemize}
  \item The parser adds these relations into the Node Model as:
  \begin{itemize}[leftmargin=1.5em]
    \item \texttt{supported\_by(parent, child)}
    \item \texttt{in\_context\_of(parent, context)}
  \end{itemize}
\end{itemize}

Dangling references (IDs mentioned here that have no corresponding node header) are treated as \textbf{errors} (see Section~\ref{sec:errors}).

\bigskip
\hrule
\bigskip

\section{Style and Clarity Guidelines}

\subsection{Goal and Module Bodies}

To keep the ACO both human-understandable and machine-interpretable:
\begin{itemize}[leftmargin=2em]
  \item Goal and Module bodies should state \textbf{claims}, not vague aspirations.
  \item Avoid ambiguous terms such as:
  \begin{itemize}[leftmargin=1.5em]
    \item ``may'', ``might'', ``possibly'', ``etc.'', ``various'', ``as appropriate''.
  \end{itemize}
  \item Prefer measurable, testable statements, or at least well-scoped qualitative claims.
\end{itemize}

\subsection{Placeholders}

Placeholders are written in \textbf{braces}:
\begin{itemize}[leftmargin=2em]
  \item \texttt{\{component\}}, \texttt{\{local policy\}}, \texttt{\{properties\}}, \texttt{\{environment\}}, etc.
\end{itemize}

Rules:
\begin{itemize}[leftmargin=2em]
  \item Placeholders should be:
  \begin{itemize}[leftmargin=1.5em]
    \item defined nearby, or
    \item conventional/common in the domain, or
    \item explained in a suitable Context node.
  \end{itemize}
  \item Tools should emit \textbf{warnings} when:
  \begin{itemize}[leftmargin=1.5em]
    \item placeholders appear that are never defined or are obviously inconsistent.
  \end{itemize}
\end{itemize}

\subsection{Naming Rules}

\begin{itemize}[leftmargin=2em]
  \item IDs (when used) must be \textbf{unique} within a single outline.
  \begin{itemize}[leftmargin=1.5em]
    \item Recommended: prefix by type (e.g.\ \texttt{G0}, \texttt{G1}, \texttt{C0}, \texttt{M2}, etc.).
  \end{itemize}

  \item \texttt{shortLabel} should:
  \begin{itemize}[leftmargin=1.5em]
    \item contain no spaces,
    \item preferably use camelCase or snake\_case.
  \end{itemize}

  \item The combination of Type, optional ID, and label should make each node recognisable.
\end{itemize}

\subsection{Text Blocks}

\begin{itemize}[leftmargin=2em]
  \item All node body text must be indented beneath the header, for example:
\end{itemize}

\begin{verbatim}
Goal G1 componentBehaviour:
  This is body text for G1.
  It may span multiple lines.

Module M1 component:
  Behaviour of the {component} meets its local policy.
\end{verbatim}

\begin{itemize}[leftmargin=2em]
  \item Blank lines may be inserted for clarity, as long as they do not break the logical grouping of body text with its header.
\end{itemize}

\bigskip
\hrule
\bigskip

\section{Error, Warning, and Undeveloped Checks (Front-End)}
\label{sec:errors}

A front-end processor parses ACO into the Node Model and performs consistency checks.

\subsection{Errors (Parsing / Structural)}

The following conditions are treated as \textbf{errors}:

\begin{itemize}[leftmargin=2em]
  \item \textbf{Duplicate node IDs}:
  \begin{itemize}[leftmargin=1.5em]
    \item The same ID used in more than one header line.
  \end{itemize}

  \item \textbf{Dangling references}:
  \begin{itemize}[leftmargin=1.5em]
    \item IDs used in explicit relation sentences that do not correspond to any declared node.
  \end{itemize}

  \item \textbf{Indentation jumps}:
  \begin{itemize}[leftmargin=1.5em]
    \item For example, going from level 0 directly to level 3 with no node at level 1 or 2.
  \end{itemize}

  \item \textbf{Unknown node types}:
  \begin{itemize}[leftmargin=1.5em]
    \item A header whose first token is not a recognised Type (Goal, Strategy, Context, Assumption, Justification, Module).
  \end{itemize}

  \item \textbf{Malformed headers}:
  \begin{itemize}[leftmargin=1.5em]
    \item Header lines missing the trailing colon (\texttt{:}),
    \item or with an ambiguous token sequence that cannot be parsed according to the rules in Section~2.2.
  \end{itemize}
\end{itemize}

Such errors should prevent further processing until fixed.

\subsection{Warnings and Undeveloped Nodes}

The following conditions are at least \textbf{warnings}:

\begin{itemize}[leftmargin=2em]
  \item Ambiguous wording in Goal/Module bodies (use of ``may'', ``etc.'', etc.).
  \item Undefined placeholders in braces \texttt{\{\dots\}}.
  \item Goals that:
  \begin{itemize}[leftmargin=1.5em]
    \item have neither supporting children (by indentation) \textbf{nor}
    \item explicit supporting evidence or modules.
  \end{itemize}
\end{itemize}

Goals with no supporting structure are considered \textbf{Undeveloped}.

Handling of Undeveloped:

\begin{itemize}[leftmargin=2em]
  \item Tools should:
  \begin{itemize}[leftmargin=1.5em]
    \item treat these as a distinct category (``Undeveloped goals/modules''), and
    \item optionally also report them as warnings.
  \end{itemize}
  \item A report may, for example:
  \begin{itemize}[leftmargin=1.5em]
    \item print a dedicated \textbf{Undeveloped} list for fast review, and
    \item also include them in the warnings list.
  \end{itemize}
\end{itemize}

This makes it easy for authors to see where an outline still lacks structure or evidence.

\bigskip
\hrule
\bigskip

\section{Worked Example --- Operational Plane Pattern}

Below is a conformant example using only indentation-defined structure, plus some optional explicit relation sentences.

\subsection{Outline Example}

\begin{verbatim}
Goal G0 operationalPlane:
  The operational plane guarantees the {local policy} is met.

  Context C0 planeDefinition:
    The system model describes the plane and its {properties}.

  Goal G1 componentBehaviour:
    Compositional behaviour of separate components ensures the local policy is met.

    Module M1 component:
      Behaviour of the {component} meets its local policy.

  Goal G2 compositionBehaviour:
    Compositional behaviour of compositions ensures the local policy is met.

    Module M2 composition:
      Behaviour of the {composition} meets the local policy.

  Module M0 interface:
    Interaction between components and compositions is as defined by the security
    architecture interface.

G0 is supported by G1, G2, M0.
G0 is in context of C0.
G1 is supported by M1.
G2 is supported by M2.
\end{verbatim}

\subsection{Derived Node Model (Semantic Graph)}

From this outline, the Node Model contains:

\textbf{Nodes} (with type, ID, and label):
\begin{verbatim}
(G0, goal,    operationalPlane)
(C0, context, planeDefinition)
(G1, goal,    componentBehaviour)
(G2, goal,    compositionBehaviour)
(M1, module,  component)
(M2, module,  composition)
(M0, module,  interface)
\end{verbatim}

\textbf{Relationships} (edges):
\begin{verbatim}
supported_by(G0, G1)
supported_by(G0, G2)
supported_by(G0, M0)
in_context_of(G0, C0)
supported_by(G1, M1)
supported_by(G2, M2)
\end{verbatim}

The Node Model is \textbf{the internal semantic representation} on which all further transformations (APL, diagrams, analysis) are based.

\bigskip
\hrule
\bigskip

\section{Mapping to APL}

The ACO notation itself is platform-neutral, but it is designed to be systematically mapped to an APL-style representation used by tooling such as O-ETB.

\subsection{General Mapping Principles}

From the Node Model:
\begin{itemize}[leftmargin=2em]
  \item Each node becomes an APL \textbf{block} with:
  \begin{itemize}[leftmargin=1.5em]
    \item a block identifier (derived from the ID or a canonicalised label),
    \item a block type (goal, strategy, context, assumption, justification, module),
    \item a label and body text.
  \end{itemize}

  \item Each \texttt{supported\_by} edge becomes a corresponding \textbf{support} relation between blocks.

  \item Each \texttt{in\_context\_of} edge becomes a corresponding \textbf{context} relation between blocks.
\end{itemize}

The precise surface syntax of APL is left to the implementation, but the mapping must be:
\begin{itemize}[leftmargin=2em]
  \item Deterministic: the same ACO outline always yields the same APL representation.
  \item Information-preserving: all node text and relations present in ACO must be recoverable from APL.
\end{itemize}

\subsection{Example (Conceptual)}

Using a Prolog-like notation for illustration (not a normative APL definition), the example in Section~7.1 could be represented conceptually as:

\begin{verbatim}
block(g0, goal,    operationalPlane,
      "The operational plane guarantees the {local policy} is met.").
block(c0, context, planeDefinition,
      "The system model describes the plane and its {properties}.").
block(g1, goal,    componentBehaviour,
      "Compositional behaviour of separate components ensures the local
       policy is met.").
block(g2, goal,    compositionBehaviour,
      "Compositional behaviour of compositions ensures the local policy
       is met.").
block(m1, module,  component,
      "Behaviour of the {component} meets its local policy.").
block(m2, module,  composition,
      "Behaviour of the {composition} meets the local policy.").
block(m0, module,  interface,
      "Interaction between components and compositions is as defined by
       the security architecture interface.").

supported_by(g0, g1).
supported_by(g0, g2).
supported_by(g0, m0).
in_context_of(g0, c0).
supported_by(g1, m1).
supported_by(g2, m2).
\end{verbatim}

Implementations are free to use different concrete syntax for APL blocks, as long as this information is faithfully represented.

\bigskip
\hrule
\bigskip

\section{Mapping to DOT (Graphviz)}

The Node Model also maps naturally to DOT/Graphviz for visualisation.

\subsection{Example DOT Rendering}

For the example in Section~7.1, a DOT rendering might be:

\begin{verbatim}
digraph OperationalPlane {
  node [shape=box, style=rounded];

  G0 [label="Goal G0: operationalPlane"];
  C0 [label="Context C0: planeDefinition",
      shape=note, style="rounded,dashed"];
  G1 [label="Goal G1: componentBehaviour"];
  G2 [label="Goal G2: compositionBehaviour"];
  M1 [label="Module M1: component"];
  M2 [label="Module M2: composition"];
  M0 [label="Module M0: interface"];

  G0 -> G1;
  G0 -> G2;
  G0 -> M0;
  G0 -> C0 [style=dashed];
  G1 -> M1;
  G2 -> M2;
}
\end{verbatim}

Styling (shapes, colours, fonts) is not part of the ACO spec, but visual conventions can help distinguish types (e.g.\ dashed edges for context).

\bigskip
\hrule
\bigskip

\section{Future Extensions}

Potential future extensions to this specification include:

\begin{itemize}[leftmargin=2em]
  \item Pattern roles (meta / property / domain / compliance patterns).
  \item Parameterised modules and reusable pattern libraries.
  \item More explicit evidence attachment constructs.
  \item Optional, richer Strategy nodes (e.g.\ AND/OR decomposition annotations).
  \item Additional relation types (e.g.\ conflict, refinement, equivalence).
\end{itemize}

These extensions should preserve:
\begin{itemize}[leftmargin=2em]
  \item the basic indentation-driven structure,
  \item the simple header syntax of \texttt{\textless Type\textgreater{} [ID] \textless shortLabel\textgreater:},
  \item and the centrality of the Node Model as the semantic representation.
\end{itemize}

\bigskip
\hrule
\bigskip

\section{Summary}

This ACO v1.0 specification establishes:

\begin{itemize}[leftmargin=2em]
  \item A simple, readable indentation-based outline format for assurance cases.
  \item A consistent header syntax that merges type, optional ID, and label.
  \item Optional natural-language relation statements for non-hierarchical links.
  \item Automated parsing from ACO into an internal Node Model.
  \item Systematic mappings from the Node Model to:
  \begin{itemize}[leftmargin=1.5em]
    \item APL-style assurance case blocks, and
    \item graphical representations such as DOT/Graphviz.
  \end{itemize}
  \item Direct generation of APL and graphical representations from a single textual source.
\end{itemize}

It is intended to evolve as MedSecurance and related work mature, and as tool integration and pattern libraries become richer.

\end{document}
